\section{Method}
\label{sec:method}

\subsection{Sub-agent harness}
\label{sec:harness}

An orchestrating Claude Code session launches each sub-agent as its own
interactive Claude Code session in a cmux terminal pane, through a small
helper script. A \texttt{launch} call prints a handle; a \texttt{collect}
call waits for the reply and closes the pane. The first agent opens one
column beside the orchestrator and later agents stack inside it. Each
session gets a fixed session ID, and we read the reply from the last
turn-end entry of the session transcript rather than from the screen, so
long replies come back whole. The prompt reaches the pane through a
temporary file, because cmux drops typed input beyond about 1,024
characters. For this study we added options for the start folder, the
permission mode (the helper refuses \texttt{bypassPermissions}) and a
settings file holding allow and deny rules.

Figure~\ref{fig:pipeline} shows one coding run. Each run gets a fresh
git worktree of the target repository at the task's base commit, with
dependencies synced, in one of eight reusable slot folders, so up to eight
runs proceed at once. Agents run in \texttt{dontAsk} permission mode with
an allowlist: running tests, read-only git, search commands, and file
edits inside the slot folders. Commits, pushes, branch switches, resets,
stashes, sub-agents and web access are denied. A call outside the
allowlist is denied without stopping the run. Claude Code asks whether to
trust any new folder and does not inherit trust from parent folders;
accepting the prompt inside a worktree saves trust on the main repository,
which then covers every worktree of it. The user accepted it once by hand.
Agents loaded the user's global Claude Code instructions, as they would in
ordinary use. Grading happens after the agent exits (Section~\ref{sec:grading}).

\begin{figure}[H]
\centering
\begin{tikzpicture}[
  node distance=0.75cm and 0.8cm,
  box/.style={rectangle, rounded corners=4pt, draw=arxivblue, thick, fill=arxivblue!8,
              font=\small, align=center, text width=2.75cm, minimum height=1.25cm,
              inner sep=4pt},
  gbox/.style={box, draw=arxivgray, fill=arxivgray!8},
  arr/.style={-{Stealth[length=2.2mm]}, thick, arxivgray}
]
% first row
\node[box] (orch) {Orchestrator session};
\node[box, right=of orch] (launch) {Launch sub-agent in a cmux pane};
\node[gbox, right=of launch] (tree) {Fresh git worktree in a trusted slot};
\node[box, right=of tree] (run) {Agent runs (\texttt{dontAsk} mode, allowlist)};
% second row, right to left
\node[gbox, below=of run] (collect) {Collect reply from transcript};
\node[box, left=of collect] (grade) {Grade from the worktree: restore tests, add hidden tests, run suite};
\node[gbox, left=of grade] (rec) {Results record};

\draw[arr] (orch) -- (launch);
\draw[arr] (launch) -- (tree);
\draw[arr] (tree) -- (run);
\draw[arr] (run) -- (collect);
\draw[arr] (collect) -- (grade);
\draw[arr] (grade) -- (rec);
\end{tikzpicture}

\caption{One coding run. The orchestrator launches the sub-agent in a
cmux pane inside a fresh worktree. After the agent exits, the reply is
read from the transcript and the run is graded from the worktree, not
from the reply.}
\label{fig:pipeline}
\end{figure}

\subsection{Models and configurations}
\label{sec:configs}

We test Claude Haiku~5.5 and Claude Sonnet~5.5 (hereafter Haiku and
Sonnet) at five effort levels: \texttt{low}, \texttt{medium},
\texttt{high}, \texttt{xhigh} and \texttt{max}. We call one model at one
effort level a \emph{config}. No Opus sub-agents were used. Not every
config ran in every experiment (Table~\ref{tab:configs}). The study ran on
November 16--17, 2026.

\begin{table}[H]
\centering
\caption{Experiments, configs and trials per cell.}
\label{tab:configs}
\small
\setlength{\tabcolsep}{6pt}
\renewcommand{\arraystretch}{1.2}
\begin{tabular}{L{4.6cm} L{6.4cm} L{2.6cm}}
\toprule
\textbf{Experiment} & \textbf{Configs} & \textbf{Trials per cell} \\
\midrule
1. Reasoning without tools & Both models at \texttt{low} to \texttt{xhigh} (8) & 1 \\
2. Coding, easy set & Both models at all five levels (10) & 1 \\
3. Coding, hard set & Haiku \texttt{low}; Sonnet \texttt{low}, \texttt{medium}, \texttt{high}, \texttt{xhigh} (5) & 3 \\
\bottomrule
\end{tabular}
\end{table}

\subsection{Tasks}
\label{sec:tasks}

Experiment~1 uses three reasoning families with exact answers and ten
difficulty levels each. Tools were forbidden, and we checked every
transcript: only thinking and text blocks appeared. Levels 6--10 ran as
per-family chains that stopped after two consecutive failures, counting a
failure at level 5; so Haiku \texttt{low}'s trace levels above 250
statements were never run.

The coding target is a private Python project of about 150 files (an agent
runtime with a daemon, several transports, model adapters, tools and a
SQLite store), with a suite of about 140 tests that runs in about 12
seconds. Its code is not published, and we describe its components only
generically. Experiments~2 and~3 use 22 tasks of four kinds, 12 in the easy
set and 10 in the hard set (Table~\ref{tab:task-sets}). Feature tasks
replay real commits in the manner of SWE-bench~\cite{jimenez2024swebench}:
the agent starts at the commit's parent, receives a written ticket (the
brief), and is graded by the commit's own tests. The two terse hard-set
features reuse two easy-set commits with much shorter briefs.

\begin{table}[H]
\centering
\caption{Task kinds and reasoning families. Coding counts are easy set plus hard set.}
\label{tab:task-sets}
\small
\setlength{\tabcolsep}{6pt}
\renewcommand{\arraystretch}{1.2}
\begin{tabular}{L{2.5cm} L{1.5cm} L{9.5cm}}
\toprule
\textbf{Kind} & \textbf{Count} & \textbf{What makes it hard} \\
\midrule
Locate & 3 & Same-named decoys; a search across Python and HTML \\
Fix & 2 + 3 & A bug two layers below its symptom; hidden tests far from the cause; two causes behind one symptom with only one visible failing test; a bug report that blames the wrong component \\
Feature & 4 + 5 & Real commits from 1 to 15 files and 25 to 602 lines added (up to 32 hidden tests); terse briefs of 70 and 115 words \\
Review & 3 + 2 & Real diffs with 0 to 3 planted bugs; the subtle ones (missing awaits, error handling, edge cases) pass the existing tests \\
\midrule
Multiplication & 10 levels & Exact long arithmetic; two numbers of 4 to 40 digits each \\
Code trace & 10 levels & Mentally running generated Python over 4 variables, 10 to 1,300 statements \\
Seating logic & 10 levels & A minimal clue set with exactly one solution; 4 to 13 people \\
\bottomrule
\end{tabular}
\end{table}

\subsection{Task construction and validation}
\label{sec:validation}

A feature task is admitted only if the commit's new tests fail before the
change and pass after it, and every other test passes both times. Sonnet
\texttt{high} sub-agents wrote the briefs, planted bugs, bug reports and
review answer keys. Planting a known fault in otherwise working code and
checking that the checks catch it follows the logic of mutation
testing~\cite{demillo1978hints}; here the checks are the visible tests for
fixes and the reviewing agent for reviews.

Because the writer also knows the answer, we verified each brief
adversarially. A Sonnet \texttt{xhigh} verifier read each brief line by
line against the hidden tests and fixed what it found. In the easy set, 2
of 4 briefs passed unchanged and 2 leaked internal helper names, which we
removed. In the hard set, 3 of 5 passed unchanged; one gained 2 missing
interface details and one gained 4, including a return type that a test
asserts.

We rejected and rebuilt one hard task, because its writer had re-planted
the easy set's bug and named the buggy function in the report.

Finally, an oracle self-test ran on all 22 tasks: the known-correct answer
must pass and doing nothing must fail. It caught one bad task, a feature
whose commit also added a dependency, so the agent's environment could not
run the new tests. The fix is that feature tasks now start with the
commit's dependency changes applied; we applied it before any run of the
two affected feature tasks. No person reviewed the answer keys; validity
rests on the oracle self-test and on spot reads by the orchestrating model.

\subsection{Grading}
\label{sec:grading}

Grading runs after the agent exits, from the worktree, not from the
agent's reply.

\begin{itemize}
  \item \textbf{Fix and feature.} Restore the original tests, add the
        hidden tests, and run the full suite, excluding two test files that
        need live network services. A run passes only with no failures and
        every hidden must-pass test passing.
  \item \textbf{Locate.} The reply must list exactly the expected set of
        \texttt{file:line} locations.
  \item \textbf{Review.} A planted bug counts as found if the reply cites
        its file within 3 lines. Extra reports on a diff with planted bugs
        are recorded but do not fail the run; a clean diff must come back
        with no reports. We adopted this rule after the pilot, because most
        pilot ``false alarms'' were real defects already present in the
        original commits, and regraded the pilot from saved replies.
  \item \textbf{Reasoning.} Exact match to the known answer.
\end{itemize}

For each run we record status, tokens by type, tool calls, denied tool
calls, agent time, and whether the agent claimed that the tests passed.
Runs that failed for harness reasons were rerun; the original row stays in
the data, marked superseded, and the rerun is counted.

\subsection{Cost accounting}
\label{sec:cost-accounting}

Dollars are API-equivalent: what the tokens would cost at first-party API
prices of November 2026, not subscription usage and not billing data. We
compute them per request from each session transcript, from token counts by
type. Sonnet costs \$2.00 per million input tokens, \$10.00 per million
output tokens and \$0.10 per million cache-read tokens. Haiku costs
\$0.10 input and \$0.50 output per million tokens for prompts up to 100K
tokens, and \$0.50 and \$2.50 beyond that; we apply this tier per request.

Two prices are assumptions rather than listed rates: cache writes at
1.25$\times$ the input rate, and Haiku cache reads at 0.1$\times$ its
input rate. A different multiplier would shift every dollar figure, so
ratios between configs are more reliable than absolute dollars. The
transcripts are not published, so the totals cannot be re-derived from the
repository alone; the method is documented in \texttt{analysis/cost\_model.py}.
We report dollars per success, meaning a config's total dollars
divided by its passed tasks.
